\chapter{Alcance}

El alcance del proyecto se define en funci\'on de un MVP acad\'emico orientado al an\'alisis asistido de resonancias magn\'eticas lumbares sagitales. El sistema propuesto no busca reemplazar el criterio profesional ni emitir diagn\'osticos, sino automatizar tareas operativas espec\'ificas vinculadas con segmentaci\'on, medici\'on, visualizaci\'on y documentaci\'on estructurada.

\section{Funcionalidades incluidas}

El proyecto incluye el dise\~no y desarrollo de un prototipo funcional con las siguientes capacidades:

\begin{itemize}
    \item Carga de estudios de RM lumbar sagital provenientes de datasets p\'ublicos o fuentes compatibles con el entorno de desarrollo.
    \item Normalizaci\'on y preprocesamiento de las im\'agenes para su uso dentro del pipeline.
    \item Segmentaci\'on autom\'atica de un conjunto acotado de estructuras anat\'omicas.
    \item Visualizaci\'on de las m\'ascaras generadas sobre la imagen original.
    \item Extracci\'on de mediciones geom\'etricas simples derivadas de las segmentaciones.
    \item Organizaci\'on de resultados por estructura anat\'omica y, cuando sea posible, por nivel lumbar.
    \item Generaci\'on de una salida estructurada editable para revisi\'on profesional.
    \item Evaluaci\'on t\'ecnica de las segmentaciones mediante m\'etricas cuantitativas.
    \item Revisi\'on cualitativa complementaria por parte de un profesional del \'area.
\end{itemize}

\section{Estructuras anat\'omicas contempladas}

El MVP se concentrar\'a en tres estructuras anat\'omicas principales:

\begin{itemize}
    \item cuerpos vertebrales;
    \item discos intervertebrales;
    \item canal espinal.
\end{itemize}

Estas estructuras se seleccionan porque se encuentran alineadas con el dataset p\'ublico definido como base del proyecto y porque permiten derivar mediciones geom\'etricas simples que pueden ser revisadas visualmente.

\section{Mediciones contempladas}

Las mediciones incluidas ser\'an acotadas y de naturaleza geom\'etrica. Podr\'an incluir, seg\'un la calidad de las m\'ascaras obtenidas y la informaci\'on espacial disponible:

\begin{itemize}
    \item \'areas proyectadas de estructuras segmentadas;
    \item dimensiones relativas de discos intervertebrales;
    \item medidas aproximadas asociadas al canal espinal en el plano sagital;
    \item relaciones o comparaciones simples entre niveles.
\end{itemize}

Estas mediciones no ser\'an interpretadas autom\'aticamente como diagn\'osticos. Por ejemplo, una reducci\'on de dimensi\'on o \'area no ser\'a clasificada por el sistema como patolog\'ia, sino presentada como dato cuantitativo revisable.

\section{Validaci\'on incluida}

La validaci\'on se desarrollar\'a en dos niveles.

En primer lugar, se realizar\'a una validaci\'on t\'ecnica cuantitativa, comparando las m\'ascaras generadas por el sistema contra anotaciones de referencia disponibles en el dataset utilizado. Para ello podr\'an emplearse m\'etricas como Dice, IoU, distancia de Hausdorff o m\'etricas de superficie, seg\'un corresponda.

En segundo lugar, se realizar\'a una validaci\'on cualitativa complementaria con un profesional vinculado al an\'alisis de im\'agenes m\'edicas. Esta revisi\'on se orientar\'a a evaluar la plausibilidad anat\'omica de las segmentaciones, la utilidad de las mediciones, la claridad de la visualizaci\'on y la legibilidad de la salida estructurada editable.

\section{Exclusiones del alcance}

Quedan fuera del alcance del proyecto:

\begin{itemize}
    \item diagn\'ostico cl\'inico autom\'atico;
    \item clasificaci\'on integral de patolog\'ias degenerativas;
    \item recomendaci\'on terap\'eutica;
    \item reemplazo del criterio m\'edico;
    \item uso cl\'inico real del sistema;
    \item validaci\'on cl\'inica multic\'entrica;
    \item integraci\'on completa con sistemas hospitalarios;
    \item integraci\'on PACS/RIS obligatoria;
    \item certificaci\'on regulatoria como software m\'edico;
    \item comparaci\'on longitudinal obligatoria entre estudios de un mismo paciente;
    \item entrenamiento sobre datos privados no anonimizados;
    \item desarrollo de un producto comercial final.
\end{itemize}

Estas exclusiones permiten mantener el proyecto dentro de un alcance acad\'emico realista y t\'ecnicamente verificable.

\section{Alcance del documento final}

El documento final abarcar\'a la fundamentaci\'on del problema, estado del arte, marco te\'orico, user research, requerimientos, arquitectura, metodolog\'ia de desarrollo, cronograma, implementaci\'on, validaci\'on, resultados, discusi\'on, conclusiones y trabajos futuros. La documentaci\'on buscar\'a dejar trazabilidad entre las decisiones de dise\~no y los resultados obtenidos durante el desarrollo del prototipo.
